% Part: computability
% Chapter: computability-theory
% Section: normal-form

\documentclass[../../../include/open-logic-section]{subfiles}

\begin{document}

\olfileid{cmp}{thy}{nfm}
\olsection{قضیهٔ صورت نرمال}

فرض کنید بتوانیم تعریف‌های توابع محاسبه‌پذیر را وصف کنیم و بیازماییم
که آیا وصفی ادعایی از سابقهٔ کامل محاسبهٔ آن تابع روی ورودی‌ای معین
درست است یا نه. آنگاه معقول است که مستقل از مدل محاسبه بتوانیم ارزش هر
تابع محاسبه‌پذیر~$f$ را روی هر ورودی~$x$ به شیوهٔ زیر تعیین کنیم:
\begin{enumerate}
  \item همهٔ وصف‌های ممکن از سوابق محاسبه را جست‌وجو کنیم.
  \item بیازماییم که آیا سابقهٔ داده‌شده، سابقهٔ محاسبهٔ~$f(x)$ است.
  \item اگر سابقهٔ درست را یافته‌ایم، ارزش~$f(x)$ را از آن استخراج
  کنیم.
\end{enumerate}
اینکه این روش واقعاً درست است محتوای قضیهٔ صورت نرمال کلینی است.

\begin{thm}[قضیهٔ صورت نرمال کلینی]
\ollabel{thm:normal-form}
رابطه‌ای بازگشتی اولیه چون~$T(e,x,s)$ و تابعی بازگشتی اولیه چون~$U(s)$
با ویژگی زیر وجود دارند: اگر~$f$ هر تابع محاسبه‌پذیر جزئی باشد، آنگاه
برای یک~$e$ و هر~$x$،
\[
f(x) \simeq U(\umin{s}{T(e, x, s)}).
\]
\end{thm}

\begin{proof}[طرح اثبات]
برای هر مدل محاسبه می‌توان وصفی از تابع محاسبه‌پذیر~$f$ را به‌دقت
تعریف کرد و آن وصف را با عدد طبیعی~$e$ کدگذاری کرد. همچنین می‌توان
مفهومی دقیق از «دنبالهٔ محاسبه» تعریف کرد که فرایند محاسبهٔ تابع با
نمایهٔ~$e$ روی ورودی~$x$ را ثبت می‌کند. چنین دنبالهٔ محاسبه‌ای را نیز
می‌توان با عددی چون~$s$ کدگذاری کرد. این کار را می‌توان چنان انجام داد
که
\begin{enumerate}
  \item رابطهٔ~$T(e,x,s)$، که هرگاه و تنها هرگاه برقرار است که عدد~$s$
  دنبالهٔ محاسبهٔ تابع با نمایهٔ~$e$ روی ورودی~$x$ را کد کند، و
  \item تابع~$U(s)$، که دنبالهٔ محاسبه با کد~$s$ را به نتیجهٔ نهایی
  آن محاسبه می‌نگارد،
\end{enumerate}
هر دو محاسبه‌پذیر باشند. در واقع رابطهٔ~$T$ و تابع~$U$ بازگشتی
اولیه‌اند.
\end{proof}

\begin{explain}
برای ارائهٔ اثباتی کامل و دقیق از قضیهٔ صورت نرمال باید مدلی از محاسبه
را تثبیت کنیم، کدگذاری وصف‌های توابع محاسبه‌پذیر و دنباله‌های محاسبه را
با جزئیات انجام دهیم و وارسی کنیم که رابطهٔ~$T$ و تابع~$U$ بازگشتی
اولیه‌اند. برای بیشتر کاربردها کافی است~$T$ و~$U$ محاسبه‌پذیر و~$U$
تام باشد.

شاید بهترین راه به‌خاطرسپردن اثبات قضیهٔ صورت نرمال این شعار باشد:
$\umin{s}{T(e,x,s)}$ دنبالهٔ محاسبهٔ تابع با نمایهٔ~$e$ را روی
ورودی~$x$ جست‌وجو می‌کند و، اگر چنین دنباله‌ای یافت شود، $U$ خروجی آن
را بازمی‌گرداند.
\end{explain}

اگر مدل محاسبه همان توابع بازگشتی جزئی باشد (مدلی که کلینی در اصل به
کار برد)، این قضیه نشان می‌دهد برای تعریف هر تابع تنها یک بار استفاده
از جست‌وجوی نامحدود، یعنی تنها یک عملگر
$\umin{y}{f(\vec x,y)}$، لازم است. بدین معنا، قضیه نشان می‌دهد هر تابع
بازگشتی جزئی صورتی نرمال دارد.

می‌توان از~$T$ و~$U$ برای تعریف برشماری
$\cfind{0}$، $\cfind{1}$، $\cfind{2}$، \dots استفاده کرد. ازاین‌پس
فرض می‌کنیم انتخاب مناسبی از~$T$ و~$U$ را تثبیت کرده‌ایم و معادلهٔ
\[
\cfind{e}(x) \simeq U(\umin{s}{T(e,x,s)})
\]
را 
\emph{تعریفِ}~$\cfind{e}$ می‌گیریم.

این هم واقعیت سودمند دیگری:

\begin{thm}
هر تابع محاسبه‌پذیر جزئی بی‌نهایت نمایه دارد.
\end{thm}

این نیز از نظر شهودی روشن است. از هر وصف یک تابع محاسبه‌پذیر می‌توان
وصف متفاوتی ساخت که همان تابع (همان زوج‌های ورودی--خروجی) را محاسبه
می‌کند، اما مثلاً نخست کاری بی‌اثر بر محاسبه انجام می‌دهد (برای نمونه،
می‌آزماید که آیا $0=0$، یا تا~$5$ می‌شمارد، و جز آن). نمایهٔ وصف
تغییریافته با نمایهٔ اصلی متفاوت خواهد بود. هر دو نمایه‌های یک تابع‌اند
که تنها به شیوه‌هایی اندکی متفاوت محاسبه شده است.

\end{document}
